\section{Introduction}
\label{sec:introduction}

Aircraft final assembly involves a large number of precedence-constrained operations, long production cycles, restricted workspaces, and a strong dependence on skilled manual labor. Aircraft pulse assembly lines (APALs) organize these operations across a sequence of physical stations and move each aircraft forward at a fixed takt, thereby making station balance and resource coordination central to production efficiency~\cite{bao2023afal,ai2025,wen2026ginppo}. In practice, APAL planning also has a dynamic scope: a baseline schedule is generated before execution, while task delays may subsequently make its operation release times and resource arrangements infeasible or inefficient. The operational setting considered in this study therefore follows a predictive--reactive scheduling paradigm, in which an initial plan is revised when a disruption is observed~\cite{vieira2003,moghaddam2022}. Here, ``predictive'' refers to the advance generation of a baseline schedule; it does not imply that the occurrence or duration of a future delay is forecast.

The APAL problem studied in this paper requires a coupled decision over an operation, a physical station, and a concrete worker team. For each selected operation, the required team cardinality and skill are given, but the identities of the workers and their station assignment must be determined during scheduling. Preassigning every worker to a station is restrictive because the workload and skill demand of each station emerge only after operations are sequenced and assigned. We therefore initialize workers as unbound resources. Once a station selects a worker as part of an operation team, that worker becomes bound to the station for the remainder of the planning horizon. Consequently, station staffing is formed endogenously through sequential scheduling decisions, allowing labor capacity to follow the emerging station loads while preserving a stable station--worker relationship. This decision structure is more specific than the operation--machine selection used in mainstream learning formulations of job-shop and flexible job-shop scheduling~\cite{zhang2020,park2021,wang2024}: APAL scheduling must jointly construct operation order, physical-station allocation, and a skill-feasible set of individual workers under pulse-line constraints.

This joint decision creates difficulties that cannot be addressed by enlarging a conventional single-step action space alone. First, the state contains heterogeneous entities and relations: operation precedence, station compatibility and load, skill requirements, worker capabilities, availability, and binding status have different semantics but interact in every decision. Second, station staffing is both a consequence of earlier scheduling actions and a constraint on later actions, so operation sequencing and resource configuration cannot be optimized independently. Third, if an operation requires $q_i$ workers from $n_i$ eligible candidates, its candidate teams grow as $\binom{n_i}{q_i}$ even though $q_i$ itself is known. Finally, selecting an operation changes the feasible stations, selecting a station changes the feasible workers, and each selected worker changes the remaining team candidates. Task-delay rescheduling adds another layer of state dependence because completed or ongoing operations must be frozen, established worker--station bindings must be inherited, and revisions must be evaluated relative to the baseline schedule. A suitable policy must therefore combine heterogeneous global reasoning, structured action construction, and feasibility maintenance after every subdecision.

Existing research addresses important parts of this problem, but not their full combination. Classical assembly-line studies provide mature formulations for balancing, sequencing, and resource-constrained task assignment~\cite{becker2006,boysen2007,battaia2013,boysen2009}. More recent aircraft final-assembly studies consider resource-constrained balancing, worker movement and allocation, skilled-worker scheduling, and APAL scheduling with operator constraints~\cite{bao2023afal,wang2024worker,liu2025afalworker,wen2025qlpso}. These methods establish the importance of labor resources, but they typically solve a particular instance through mathematical programming, heuristics, or metaheuristics, and address balancing, task scheduling, or worker allocation within problem-specific decision structures. Graph neural networks and deep reinforcement learning provide a complementary route to reusable scheduling policies~\cite{cappart2023,zhang2020,park2021,wang2024,panzer2022}; direct APAL studies have also begun to explore graph-enhanced and skill-aware policy learning~\cite{ai2025,wen2026ginppo}. Nevertheless, existing learning actions mainly determine operation priorities or operation--machine/station assignments, rather than explicitly decoding a concrete fixed-cardinality worker team while station staffing is being formed. Predictive--reactive rescheduling research, in turn, explains how a baseline plan can be revised after an event~\cite{vieira2003,moghaddam2022}, but does not by itself represent the APAL-specific inheritance of frozen operations, station restrictions, and endogenous worker bindings. The remaining problem is thus not merely to add worker constraints to an existing scheduler, but to align a multi-resource APAL representation with a structured joint decoder and stepwise constraint enforcement.

To address this problem, we develop \method{}, a graph reinforcement learning framework whose policy network, \network{}, is organized into three aligned components. First, an APAL heterogeneous graph represents operations, stations, virtual skill-hub nodes, and workers, while relation-aware graph attention and global mean pooling combine local resource interactions with system-level load information. Second, a feasibility-constrained autoregressive pointer decoder selects a ready operation, a compatible station, and then exactly $q_i$ distinct workers; masks are refreshed after each subdecision to enforce precedence, spatial, capacity, skill, availability, non-duplication, and worker-binding constraints. This factorization replaces a flattened combinatorial action with a sequence of conditional choices without allowing the policy to choose the required team size. Third, proximal policy optimization (PPO) trains the shared policy~\cite{schulman2017}; size-independent node processing, global pooling, and autoregressive decoding keep the representation and action interface compatible with different instance sizes, while feature extension and warm-start training adapt the same representation--decoder backbone to task-delay rescheduling. These capabilities are examined under two deliberately separate protocols: a single locked initial-scheduling checkpoint is evaluated across instance scales without fine-tuning, whereas the rescheduling study extends the state features, initializes from the baseline policy, and continues training. This distinction prevents architectural support for variable graph sizes from being conflated with empirical cross-scale generalization, and separates direct policy reuse from warm-start transfer.

The remainder of this paper is organized as follows. Section~\ref{sec:related} reviews APAL and assembly-line optimization, learning-based scheduling, and dynamic rescheduling. Section~\ref{sec:formulation} formulates the initial APAL scheduling problem, its task-delay extension, and the corresponding Markov decision process. Section~\ref{sec:method} presents the heterogeneous representation, autoregressive decoder, and PPO learning and transfer procedure. Section~\ref{sec:experiments} describes the experimental protocols and reports the empirical results. Section~\ref{sec:discussion} discusses the findings and limitations and concludes the paper.
